% Reader copy reviewed 3 October 2026; frozen source and receipts remain unchanged.
\begin{proof}
Write $c(y)=u(\eta,y)$ near the node.  We use the comparison theorem
\Cref{thm:p071-infinite-star-collapse-candidate}, including its endpoint
trace, slope budget, finite-activity truncation, and finite harmonic
smoothing.

\medskip\noindent
\emph{Locality and decreasing closure of the comparison conditions.}
The definition of $\mathcal J_X$ is local, so restriction and gluing are
immediate once the product junction inequality is known to be local.  To
check this point, contract every point occurring in
\eqref{eq:p071-infinite-euclidean-jensen} by a common factor
$0<s<1$ towards $(\eta,y_0)$:
\[
 (-r_i,y_i)\longmapsto
 (-sr_i,\,y_0+s(y_i-y_0)).
\]
The barycentre relation is preserved, with every reciprocal weight
multiplied by $s^{-1}$.  For sufficiently small $s$, all contracted
points lie in one comparison chart at $(\eta,y_0)$.  Write the exceptional
set of that restricted chart as $A\cup E$, enlarging $E\subset B\setminus
A$ harmlessly if necessary.  Retain the contracted data on every original
tuple branch in $A\cup F$.  For each newly exceptional branch
$e\in E\setminus F$, choose any $0<\sigma_e$ in its restricted collar
and add the sample $(-\sigma_e,y_0)$.  The enlarged barycentre relation
still holds, because the new point and the centre have the same Euclidean
coordinate.  With
\[
                  \Theta=\sum_{e\in E\setminus F}\frac1{\sigma_e},
\]
the restricted-chart inequality is
\begin{align*}
 (s^{-1}\Lambda+\Theta)c(y_0)
 &\leq \sum_{b\in F}
   \frac{u_b(-sr_b,y_0+s(y_b-y_0))}{sr_b}\\
 &\quad+\sum_{a\in A}
   \frac{u_a(-s\rho_a,y_0+s(y_a-y_0))}{s\rho_a}
   +\sum_{e\in E\setminus F}
   \frac{u_e(-\sigma_e,y_0)}{\sigma_e}.
\end{align*}
Each newly exceptional branch was full before restriction, so its last
sample is at most $c(y_0)$.  Joint convexity on each original segment gives
\[
 u_i(-sr_i,y_0+s(y_i-y_0))
 \leq (1-s)c(y_0)+s u_i(-r_i,y_i).
\]
Substitution bounds the right side of the restricted-chart inequality by
\[
 (s^{-1}-1)\Lambda c(y_0)
 +\sum_{b\in F}\frac{u_b(-r_b,y_b)}{r_b}
 +\sum_{a\in A}\frac{u_a(-\rho_a,y_a)}{\rho_a}
 +\Theta c(y_0).
\]
Cancelling the two added central contributions proves the original
inequality.  The same argument proves compatibility when a full branch
becomes exceptional after restriction.  Thus the displayed conditions do
define a sheaf.

We next prove sequential closure.  Let $u_n\downarrow u$ on an open set
and fix a node point.  Choose a standard product
$U_0\times V_0$ in that open set and a compact convex set
$Q\Subset V_0$ whose interior contains the Euclidean coordinate of the
point.  For fixed $n$, compactness of $\{\eta\}\times Q$ supplies
finitely many comparison charts for $u_n$.  The union of their finite
exceptional sets and countable activity sets is countable; call it
$K_n$.  Local joint convexity and the contraction argument above extend
the branch and junction conditions across $U_0\times\operatorname{int}Q$.
Moreover,
\begin{equation}
 u_n(p_b(t),y)=u_n(\eta,y)
 \quad
 \bigl(b\notin A\cup K_n,\ t\in[-\infty,0],\ y\in Q\bigr),
 \label{eq:p071-infinite-euclidean-limit-inactive}
\end{equation}
where $A$ is the finite exceptional set of $U_0$.  Hence the limit has
one countable activity set $\bigcup_nK_n$.

Joint extended-real convexity passes to a decreasing limit.  For each fixed
$y$, the endpoint trace passes by the argument in the exact-star theorem:
monotonicity gives one inequality, and a sufficiently advanced row gives
the reverse inequality, also when the limiting endpoint value is
$-\infty$.  The same specialization, with every Euclidean sample equal
to $y$, shows that each fibre is a member of $\mathscr J_B$.
Suppose now that $c(y_0)$ is finite in a prescribed product junction
tuple.  The scalar exact-star argument shows that every finite-stratum
value in that tuple is finite.  Consequently all terms in the junction
inequality for $u_n$ converge in $\mathbf R$, and termwise passage to
the limit proves \eqref{eq:p071-infinite-euclidean-jensen} for $u$.
Thus $\mathcal J_X$ is closed under pointwise decreasing sequential
limits.

Every generator belongs to $\mathcal J_X$.  Only the product junction
condition requires proof.  Put
$d_i=\partial_tg_i(0,y_0)$.  Joint convexity and smoothness give
\[
 g_i(-r_i,y_i)\geq
 c_g(y_0)-r_id_i+
 \nabla c_g(y_0)\mathbin\cdot(y_i-y_0),
\]
where $r_i$ denotes $r_b$ or $\rho_a$.  Divide by $r_i$ and sum
over the selected full branches and all exceptional branches.  The
Euclidean linear terms cancel by
\eqref{eq:p071-infinite-euclidean-barycentre}.  Every omitted full-branch
derivative is nonnegative, since that branch is convex and constant near
$-\infty$.  The balance law
\eqref{eq:p071-infinite-euclidean-smooth-balance} therefore says that the
sum of the selected derivatives is nonpositive.  The remaining inequality
is exactly \eqref{eq:p071-infinite-euclidean-jensen}.  Minimality now gives
\begin{equation}
                         \mathcal P_X\subseteq\mathcal J_X.
 \label{eq:p071-infinite-euclidean-minimality}
\end{equation}

\medskip\noindent
\emph{Partial Legendre reduction at a finite node value.}
Let $u\in\mathcal J_X$ near $(\eta,y_0)$, and first suppose that
$c(y_0)\in\mathbf R$.  After shrinking the Euclidean ball, $c$ is
finite: an extended-real convex function on an open convex set cannot take
both a finite value and the value $-\infty$.  Specializing the product
junction condition to equal Euclidean coordinates places every fibre in
$\mathscr J_B$.  The finite-centre part of the exact-star theorem then
shows that all finite-stratum branch values are real.

Use a comparison chart on $U\times V_0$, with exceptional set $A$ and
countable activity set $K$.  Shrink the exceptional collars once.  Let
$D\subset U$ be the resulting compact closed standard core.  Choose an
open convex set $V\ni y_0$ and a compact convex set $Q$ such that
\[
                  \overline V\subset\operatorname{int}Q,
                  \qquad Q\Subset V_0.
\]
For
$p\in(\mathbf R^d)^*$, define
\begin{align}
 f_0(p)&=\min_{y\in Q}\bigl(c(y)-p\mathbin\cdot y\bigr),
 \label{eq:p071-infinite-euclidean-legendre-centre}\\
 f_b(t,p)&=\min_{y\in Q}
       \bigl(u_b(t,y)-p\mathbin\cdot y\bigr).
 \label{eq:p071-infinite-euclidean-legendre-branch}
\end{align}
At a full endpoint the second minimum is interpreted in the extended
order.  Its endpoint slice is either finite on the interior of $Q$ or
identically $-\infty$; these are the only possibilities for an
extended-real convex function on a connected open convex set.

For every fixed $p$, the functions in
\eqref{eq:p071-infinite-euclidean-legendre-centre}--%
\eqref{eq:p071-infinite-euclidean-legendre-branch} form a member
$f(\,\cdot\,,p)\in\mathscr J_B(D)$.  Convexity in the branch variable
follows by taking minimizing Euclidean points at two branch parameters and
using their convex combination in $Q$.  For a full branch, exact endpoint
trace passes through the minimum: the lower bound follows from
monotonicity, while evaluation at a minimizing endpoint point gives the
upper bound; if the endpoint slice is $-\infty$, evaluation at any fixed
Euclidean point gives convergence to $-\infty$.

It remains to check the finite-point balance.  Choose minimizing points
$y_i\in Q$ for all finitely many branch terms and set
\[
 y_*=\Lambda^{-1}
 \left(\sum_{b\in F}\frac{y_b}{r_b}
       +\sum_{a\in A}\frac{y_a}{\rho_a}\right)\in Q.
\]
Then the product junction inequality gives
\begin{align*}
 \Lambda f_0(p)
 &\leq \Lambda\bigl(c(y_*)-p\mathbin\cdot y_*\bigr)\\
 &\leq
 \sum_{b\in F}\frac{f_b(-r_b,p)}{r_b}
 +\sum_{a\in A}\frac{f_a(-\rho_a,p)}{\rho_a}.
\end{align*}
This is the Jensen form of the scalar exact-star balance.  Finally,
\eqref{eq:p071-infinite-euclidean-countable-activity} gives
\begin{equation}
 f_b(t,p)=f_0(p)
 \quad(b\in B\setminus(A\cup K)).
 \label{eq:p071-infinite-euclidean-legendre-inactive}
\end{equation}

Conversely, suppose that scalar profiles $h_b(t,p)$ belong to
$\mathscr J_B(D)$ for every $p$.  Then
\begin{equation}
 H_b(t,y)=\sup_p\bigl(p\mathbin\cdot y+h_b(t,p)\bigr)
 \label{eq:p071-infinite-euclidean-reconstruction}
\end{equation}
is jointly convex on every branch and satisfies the product junction
condition.  Indeed, multiply the scalar Jensen inequality by the
reciprocal-distance weights, add the Euclidean barycentre identity, and
then take the supremum in $p$.  This observation is the bridge between
the scalar exact-star smoothing and the mixed Hessian required here.

\medskip\noindent
\emph{Finite activity and finite dual sampling.}
Choose nested finite sets $K_n\subset K$ with union $K$, stopping if
$K$ is finite.  Define scalar profiles, for every $p$, by
\begin{align}
 h_{n,a}(t,p)&=\max\{f_a(t,p),f_0(p)-n\}, &&a\in A,
 \label{eq:p071-infinite-euclidean-truncate-a}\\
 h_{n,b}(t,p)&=\max\{f_b(t,p),f_0(p)-n\}, &&b\in K_n,
 \label{eq:p071-infinite-euclidean-truncate-k}\\
 h_{n,b}(t,p)&=f_0(p), &&b\in B\setminus(A\cup K_n).
 \label{eq:p071-infinite-euclidean-truncate-rest}
\end{align}
The slope-budget and truncation argument in the exact-star theorem shows
that each $h_n(\,\cdot\,,p)$ is a finite continuous, finite-active member
of $\mathscr J_B(D)$.  It also gives
\begin{equation}
 h_{n+1}(x,p)\leq h_n(x,p),
 \qquad h_n(x,p)\longrightarrow f_x(p).
 \label{eq:p071-infinite-euclidean-profile-limit}
\end{equation}
On a newly retained branch, the first inequality uses
$f_b(t,p)\leq f_0(p)$, the monotonicity of a full scalar comparison
profile towards the node.

Let $H_n$ be reconstructed from $h_n$ by
\eqref{eq:p071-infinite-euclidean-reconstruction}.  Local convex
biconjugacy on $\operatorname{int}Q$ gives the explicit formula
\begin{equation}
 H_{n,b}(t,y)=
 \begin{cases}
  \max\{u_b(t,y),c(y)-n\},&b\in A\cup K_n,\\
  c(y),&b\in B\setminus(A\cup K_n).
 \end{cases}
 \label{eq:p071-infinite-euclidean-truncated-reconstruction}
\end{equation}
The formula remains valid at a full endpoint: if the endpoint slice is
finite, use its supporting covectors; if it is identically $-\infty$,
only the second entry of the maximum remains.  Hence
\begin{equation}
 H_{n+1}\leq H_n,
 \qquad H_n\longrightarrow u
 \quad\text{pointwise on }D\times V.
 \label{eq:p071-infinite-euclidean-H-limit}
\end{equation}

Each $H_n$ is finite and continuous on
$D\times\overline V$.  On finite branch strata this follows from finite
joint convexity.  We record why pointwise continuity in the branch
variable is enough at the node.  If a finite convex function $q$ on
$[-R,0]\times V_0$ has continuous boundary value
$q(0,y)=c_q(y)$ and
$q(t,y)\to c_q(y)$ for every fixed $y$, then $q$ is continuous at
every point of $\{0\}\times V_0$.  For the upper bound, write a nearby
$(t,y)$ as a convex combination of $(0,z)$ and $(-R,w)$, with
$w$ fixed and the coefficient of $(-R,w)$ tending to zero.  For the
lower bound, first choose $t_0<0$ so that
$q(t_0,y_0)$ is close to $c_q(y_0)$, then write
$(t_0,y_0)$ as a convex combination of the nearby $(t,y)$ and a point
$(-R,z)$.  Solving the convexity inequality for $q(t,y)$, and using
boundedness of the finite convex slice $q(-R,\,\cdot\,)$ on a compact
Euclidean neighbourhood, gives the lower bound.  This proves the claim.

Only finitely many branches of $H_n$ are active.  Apply the boundary lemma
to each of them.  The scalar endpoint-slope bounds give its fixed-$y$
hypothesis, while every inactive branch is exactly $c(y)$.  The
cofinite-branch topology now gives continuity at the node.

At a full endpoint, choose any sequence of branch parameters decreasing to
$-\infty$.  Equations
\eqref{eq:p071-infinite-euclidean-endpoint-trace} and
\eqref{eq:p071-infinite-euclidean-truncated-reconstruction} give monotone
convergence of the corresponding continuous Euclidean slices to the finite
convex endpoint slice.  Dini's theorem makes this convergence uniform on
$\overline V$; branch monotonicity then makes the estimate uniform for all
sufficiently terminal parameters.  Since $D$ is compact, so is
$D\times\overline V$.

Fix a summable positive sequence $(\varepsilon_n)$.  For each point of
that compact product, choose a covector which realizes
\eqref{eq:p071-infinite-euclidean-reconstruction} within
$\varepsilon_n$.  Continuity of the corresponding scalar profile on
$D$ supplies a neighbourhood on which the same choice works.  A finite
subcover gives covectors $p_{n,1},\ldots,p_{n,q_n}$ such that
\begin{equation}
 H_n-\varepsilon_n\leq
 A_n:=\max_j\bigl(p_{n,j}\mathbin\cdot y
                 +h_n(x,p_{n,j})\bigr)
 \leq H_n
 \quad\text{on }D\times\overline V.
 \label{eq:p071-infinite-euclidean-finite-dual-sampling}
\end{equation}
No compactness of the full dual space is used.

For each selected profile, truncation is inactive on a sufficiently short
node collar.  Its endpoint derivatives are therefore the finite derivatives
of the corresponding scalar comparison profile, and their sum is
nonpositive by the scalar slope balance.  Choose once a smaller standard
neighbourhood $U'$ with closure in $D$.  The finite harmonic smoothing
from the exact-star theorem then gives, for every selected covector, a smooth
psh scalar function $s_{n,j}$ on that same $U'$ with
\[
 \bigl\lVert s_{n,j}-h_n(\,\cdot\,,p_{n,j})\bigr\rVert_{U'}
 <\varepsilon_n.
\]
Thus $s_{n,j}+2\varepsilon_n$ lies between the scalar profile plus
$\varepsilon_n$ and that profile plus $3\varepsilon_n$.  The lifted
function
\[
 (x,y)\longmapsto
 s_{n,j}(x)+2\varepsilon_n+p_{n,j}\mathbin\cdot y
\]
is an exact product generator.  Apply a smooth convex coordinatewise
nondecreasing regularized maximum to these finitely many lifts, with error
at most $\varepsilon_n$, and add one further $\varepsilon_n$.  The
result is a generator $W_n$ on $U'\times V$ satisfying
\begin{equation}
                         H_n\leq W_n\leq H_n+5\varepsilon_n.
 \label{eq:p071-infinite-euclidean-smoothing-error}
\end{equation}
The balance law survives the regularized maximum.  At the node, its partial
derivatives with respect to the finitely many inputs are independent of the
branch, so summing first over the branches reduces
\eqref{eq:p071-infinite-euclidean-smooth-balance} to the corresponding
balance for each input.  A finite union of activity sets is finite, and a
common terminal collar exists for every full active branch.  Joint
convexity follows from convex monotone composition.

Set
\[
 d_n=5\sum_{k>n}\varepsilon_k,
 \qquad G_n=W_n+d_n.
\]
Equations \eqref{eq:p071-infinite-euclidean-H-limit} and
\eqref{eq:p071-infinite-euclidean-smoothing-error} give
\[
 G_{n+1}\leq H_{n+1}+5\varepsilon_{n+1}+d_{n+1}
 \leq H_n+d_n\leq G_n.
\]
Moreover, $G_n\to u$ pointwise.  This proves the fixed-neighbourhood
first-stage approximation at every finite node value.

\medskip\noindent
\emph{The value $-\infty$, full endpoints, and interior points.}
If $c(y_0)=-\infty$, joint convexity makes $c$ identically
$-\infty$ after a Euclidean shrinking.  On every full branch, the scalar
exact trace and convexity then force the whole branch to be
$-\infty$; on the finitely many exceptional collars, joint convexity
forces $-\infty$ after one common shrinking.  Thus the germ is
identically $-\infty$, and the constants $-n$ give the required
approximation.

It remains to treat a point $(-\infty,y_0)$ in a full branch.  Work on
$[ -\infty,R)\times V_0$, where $R<0$.  The identically $-\infty$
case is again immediate.  Otherwise the function is finite on the finite
stratum and is nondecreasing in the branch variable.  Choose $r<R_1<R$,
an open convex set $V\ni y_0$, and a compact convex set $Q$ with
$\overline V\subset\operatorname{int}Q\Subset V_0$, and put
\begin{align*}
 f(t,p)&=\min_{y\in Q}\bigl(u(t,y)-p\mathbin\cdot y\bigr),\\
 f_1(p)&=\min_{y\in Q}\bigl(u(R_1,y)-p\mathbin\cdot y\bigr),\\
 h_n(t,p)&=\max\{f(t,p),f_1(p)-n\}.
\end{align*}
Then $h_n(\,\cdot\,,p)$ is finite, convex, nondecreasing, and has the
exact endpoint trace.  Its reconstruction is
\begin{equation}
 H_n(t,y)=\max\{u(t,y),u(R_1,y)-n\},
 \label{eq:p071-infinite-euclidean-endpoint-truncation}
\end{equation}
so $H_n\downarrow u$ and every $H_n$ is continuous on the compact
terminal core $[ -\infty,r]\times\overline V$.  The endpoint continuity
again follows from Dini's theorem.

Fix a summable positive sequence $(\varepsilon_n)$.  Finite dual sampling
on the compact terminal core gives covectors
$p_{n,1},\ldots,p_{n,q_n}$ such that
\[
 H_n-\varepsilon_n\leq
 A_n:=\max_j\bigl(p_{n,j}\mathbin\cdot y
                 +h_n(t,p_{n,j})\bigr)\leq H_n.
\]
For each selected scalar profile, use the endpoint construction in the
proof of \Cref{thm:p071-infinite-star-collapse-candidate}.  It supplies a
decreasing sequence of smooth convex functions, constant on terminal
collars, with limit $h_n$.  Since $h_n$ is finite and continuous on the
compact endpoint interval, Dini's theorem permits the choice of one term
$s_{n,j}$ satisfying
\[
       h_n(\,\cdot\,,p_{n,j})\leq s_{n,j}\leq
       h_n(\,\cdot\,,p_{n,j})+\varepsilon_n.
\]
Add $2\varepsilon_n$ to the finitely many lifted functions
$s_{n,j}(t)+p_{n,j}\mathbin\cdot y$, apply the same regularized maximum
with error at most $\varepsilon_n$, and add one further
$\varepsilon_n$.  The resulting exact product generator $W_n$ satisfies
\[
                       H_n\leq W_n\leq H_n+5\varepsilon_n
\]
on one fixed terminal product neighbourhood.  The neighbourhood is
independent of $n$; only finitely many selected profiles occur at each
stage, so their constant collars have a common terminal subcollar.  Set
\[
 d_n=5\sum_{k>n}\varepsilon_k,\qquad G_n=W_n+d_n.
\]
Since $H_{n+1}\leq H_n$, the same calculation as at the node gives
\[
 G_{n+1}\leq H_{n+1}+5\varepsilon_{n+1}+d_{n+1}
 \leq H_n+d_n\leq G_n.
\]
Moreover $G_n\to u$ pointwise.  Thus the $G_n$ give the required
decreasing exact product generators on one fixed terminal neighbourhood.

At an interior branch point, a locally extended-real jointly convex
function is either identically $-\infty$ or finite.  In the finite case,
the point lies in an ordinary open convex chart in
$\mathbf R^{d+1}$.  The one-point Stone case of
\Cref{thm:p071-stone-euclidean-collapse} therefore gives a decreasing
smooth jointly convex approximation on one fixed smaller chart.  We have
therefore proved
$\mathcal J_X\subseteq\mathcal P_{X,1}$ at every point.  Together with
\eqref{eq:p071-infinite-euclidean-minimality} and the tautological inclusion
$\mathcal P_{X,1}\subseteq\mathcal P_X$, this proves
\eqref{eq:p071-infinite-euclidean-collapse} and the fixed-neighbourhood
assertions.
\end{proof}
